\Introduction

\textbf{Актуальность.}
Управление цифровым присутствием является одной из ключевых задач для субъектов малого и среднего бизнеса (далее~--- МСБ) в~условиях цифровой экономики. Информация о~компании распределена по множеству платформ~--- геосервисам, социальным сетям, мессенджерам~--- каждая из которых обладает собственным интерфейсом управления. Актуализация данных (часов работы, адреса, контактов), мониторинг отзывов и публикация контента выполняются вручную на каждой площадке, что создаёт риск рассогласования данных и требует значительных трудозатрат.

Проблема усугубляется гетерогенностью интеграционных интерфейсов целевых платформ: одни предоставляют официальные программные интерфейсы приложений (API), тогда как другие~--- включая лидирующие на российском рынке геосервисы Яндекс.Бизнес (53~\% пользователей) и 2ГИС (44~\%)~--- не имеют публичных API для программного управления карточкой организации. Это делает невозможным создание универсального решения только на базе API и требует гибридного подхода, сочетающего работу через API с~роботизированной автоматизацией процессов (RPA) для платформ без программного интерфейса.

Существующие решения~--- системы управления взаимоотношениями с~клиентами (CRM, \textit{customer relationship management}): Bitrix24, SendPulse, Kommo; инструменты планирования публикаций в~социальных сетях (SMM, \textit{social media management}): Hootsuite, Buffer, SMMplanner; платформы мониторинга репутации (ORM, \textit{online reputation management}): YouScan~--- решают отдельные подзадачи, однако ни одно из них не обеспечивает комплексного управления консистентностью данных на разнородных платформах с~единой оркестрацией действий специализированных агентов.

Мультиагентные системы (МАС) предоставляют архитектурную основу для решения указанной проблемы: специализированные агенты инкапсулируют метод интеграции с~конкретной платформой, а~центральный оркестратор координирует их работу через единый протокол. Такой подход обеспечивает расширяемость~--- добавление новых платформ не требует модификации ядра системы.

\textbf{Объект исследования}~--- процесс управления цифровым присутствием субъектов малого и~среднего бизнеса на внешних платформах.

\textbf{Предмет исследования}~--- автоматизация управления цифровым присутствием на основе мультиагентной архитектуры с~гибридной моделью интеграции (API~+ RPA).

\textbf{Цель работы}~--- разработка программного комплекса на основе мультиагентной архитектуры с~гибридной моделью интеграции для управления цифровым присутствием субъектов МСБ.

Для достижения указанной цели необходимо выполнить следующие задачи:

\begin{enumerate}
  \item провести анализ предметной области и~существующих решений;
  \item сформулировать функциональные и~нефункциональные требования к~системе;
  \item обосновать выбор технологий и~архитектурных решений;
  \item спроектировать архитектуру мультиагентной системы;
  \item спроектировать модель данных;
  \item реализовать прототип программного комплекса;
  \item провести тестирование и~оценку эффективности разработанного решения.
\end{enumerate}

\textbf{Методы исследования:} сравнительный анализ существующих решений методом взвешенных критериев; объектно-ориентированное проектирование; моделирование бизнес-процессов (AS-IS / TO-BE); инфологическое и~даталогическое моделирование данных; микросервисная декомпозиция; функциональное тестирование по критическому пути пользователя.

\textbf{Практическая значимость} работы состоит в~создании программного комплекса, позволяющего субъектам МСБ управлять цифровым присутствием на нескольких площадках из единого интерфейса с~использованием естественного языка. Система автоматизирует рутинные операции (обновление данных, публикация контента, мониторинг отзывов), снижает риск рассогласования информации между платформами и~обеспечивает журналирование всех действий. Гибридная модель интеграции (API~+ RPA) позволяет работать в~том числе с~платформами, не предоставляющими публичного API, что является критичным для российского рынка геосервисов.

\textbf{Структура работы.} Работа состоит из введения, трёх глав, заключения, списка использованных источников и~приложения.

В~\textit{первой главе} проведён анализ предметной области, выполнен сравнительный анализ существующих решений, обоснован выбор технологий и~сформулированы требования к~системе.

В~\textit{второй главе} описана архитектура системы, модель данных, проектирование пользовательского интерфейса и~средства реализации.

В~\textit{третьей главе} приведены результаты тестирования и~оценка эффективности разработанного решения.
